\section{Lipschitz representatives}
\label{app:lipschitz-representatives}

We prove the uniform replacement lemma used to construct the compact
space of quasi-isometry pairs in Section~\ref{sec:stationary}.

\begin{proof}[Proof of Lemma~\ref{lem:lipschitz-pairs}]
A coarse inverse can first be chosen with constants depending only on
$K,C$.  Indeed, for each target point choose a source point whose image
is within the coarse-surjectivity constant.  The two quasi-isometry
inequalities bound the distortion of this choice and its two composition
errors.  It is therefore enough to smooth either member of a uniformly
controlled coarse-inverse pair by a bounded amount.

Choose a maximal $1$-separated set $\mathcal V\subset G$.
Properness makes it locally finite, and maximality makes its closed
unit balls cover $G$.  The cover by open balls of radius $2$ has uniformly
bounded multiplicity: disjoint balls of radius $1/3$ about participating
centers lie in a ball of fixed radius, and left invariance gives the
same Haar-volume ratio at every center.  Its nerve therefore has bounded
dimension, say $q$; realize each simplex as a standard regular Euclidean
simplex with unit edges.
The functions
\[
 a_v(g)=\max\{0,2-d(g,v)\},\qquad
 \lambda_v(g)=\frac{a_v(g)}{\sum_{u\in\mathcal V}a_u(g)}
\]
define a partition map $\lambda$ from $G$ to the nerve.  The denominator
is at least $1$.  If $M$ bounds the cover multiplicity, normalization
and the $1$-Lipschitz cutoffs give
$\|\lambda(g)-\lambda(h)\|_1\le4M d(g,h)$.
This also controls the intrinsic piecewise Euclidean metric.  Indeed,
for probability vectors $a,b$ with $\|a-b\|_1<2$, normalize their
coordinatewise minimum to a probability vector $c$.  The vector $c$
lies in the common face of their support simplices, and the two
segments $a$ to $c$ and $c$ to $b$ have total Euclidean length at most
$2\|a-b\|_1$.  Thus $\lambda$ is locally $8M$-Lipschitz for the
intrinsic metric.  Subdivision of geodesics in $G$ makes this a global
bound.

Let $J$ be either coarse map.  Map the vertex $v$ to $J(v)$.
The image vertices of a simplex have uniformly bounded diameter, since
its source vertices have pairwise distance at most $4$.
Extend over simplices by increasing dimension, always preserving the
maps already defined on their faces.  The required controlled extension
uses exponential coordinates on $G$, which are global by
Lemma~\ref{lem:cartan-model}.  For a center $p\in G$ put
\[
 \mathcal C_p(s,z)=p\exp\bigl(s\log(p^{-1}z)\bigr),
 \qquad 0\le s\le1.
\]
For every fixed $R$, on $d(p,z)\le R$ this contraction has a uniformly
bounded derivative in $s$ and spatial Lipschitz constant at most
$c_Rs$.  To see the latter bound, differentiate the smooth map
$\exp(s\log z)$ on a compact set.  Its spatial derivative vanishes at
$s=0$, so the mean-value estimate in $s$ supplies the factor $s$.
For the spatial distance estimate, bound derivatives on an enlarged
compact ball containing geodesics between points of the original ball.
Left translation makes the bounds independent of $p$.

On a simplex write a point as $b+s(u-b)$, where $b$ is its barycenter
and $u$ lies on its boundary.  If $H$ is the previously constructed
boundary map, extend it by $\mathcal C_p(s,H(u))$, taking $p$ to be
the image of one vertex.  Facewise Lipschitz bounds give a Lipschitz
boundary map for the Euclidean distance, since the boundary of a
regular simplex of dimension at least two has bounded intrinsic
quasiconvexity constant.  For an edge its two endpoint images already
have bounded distance.  The radial extension is quantitatively
Lipschitz: if $\rho$ is the distance from $b$ to the boundary and
$D$ is the maximum distance from $b$ to the simplex, then for
$x=b+s(u-b)$ and $x'=b+s'(u'-b)$ one has
\[
 |s-s'|\le \rho^{-1}|x-x'|,\qquad
 \min(s,s')|u-u'|\le (1+D/\rho)|x-x'|.
\]
Combine these with the spatial $O(s)$ bound and the uniform time
derivative bound of $\mathcal C_p$.  Regular simplices of bounded
dimension have uniform values of the geometric constants involved.
The extension also remains in a ball of controlled radius about the
vertex images, because the contraction is evaluated on a fixed compact
ball.
Induction through at most $q$ dimensions gives one uniform Lipschitz
constant and one uniform radius.  Shared faces retain their previous
maps, so the extensions agree.  Simplexwise Lipschitz bounds give a
Lipschitz map for the intrinsic metric of the nerve.

Compose this map with $\lambda$.  Whenever $\lambda_v(g)>0$, one has
$d(g,v)<2$, hence $d(J(g),J(v))\le 2K'+C'$ for the uniform constants
of the coarse map under consideration.  The controlled radius of the
simplex extension therefore bounds the distance of the resulting
Lipschitz map from $J$, uniformly in $g$.

Apply this construction to both coarse maps.  Bounded perturbations
preserve the lower quasi-isometry inequalities.  They preserve the
bounded composition errors as well: use the original coarse upper
inequality to compare the original map at two points a bounded distance
apart, and then the bounded perturbation of its values.  Enlarging a
single constant $\Lambda$ gives all the stated bounds.
\end{proof}
